\documentclass[11pt]{article}
\usepackage[a4paper,margin=25mm]{geometry}
\usepackage[T1]{fontenc}
\usepackage{amsmath,amssymb,amsthm,mathtools}
\usepackage[hidelinks,hypertexnames=false]{hyperref}
\setlength{\emergencystretch}{2em}
\newcommand{\C}{\mathbb C}\newcommand{\R}{\mathbb R}
\newcommand{\PP}{\mathcal P}\newcommand{\Id}{\operatorname{Id}}
\newcommand{\tr}{\operatorname{tr}}\newcommand{\im}{\operatorname{im}}
\newcommand{\ord}{\operatorname{ord}}
\newtheorem{theorem}{Theorem}[section]
\newtheorem{proposition}[theorem]{Proposition}
\newtheorem{lemma}[theorem]{Lemma}
\title{The native arithmetic secular determinant:\\
complete pole cancellation and the retained quotient metric}
\author{}\date{20 September 2026}
\begin{document}\maketitle

\section{The original arithmetic source and every coordinate}

This is a finite-parameter sequel to the native polynomial calculation
\cite{Native}. Its human-source comparison is the finite rank-one
determinant calculation of Connes--Consani--Moscovici \cite{CCM},
Lemma \texttt{key} and equation \texttt{eq:detDp}.
Every determinant below is a finite determinant. No regularized
infinite determinant or convergence of spectra is inferred.
No conclusion that every root of the original arithmetic
annihilator lies on $\Re S=c$ is inferred.

Retain the actual finite divisor and source
\begin{equation}\tag{NS1}\label{NS1}
h(s)=\prod_{\rho\in Z}(s-\rho)^{m_\rho},\qquad
g(s)=2\xi(s)=s(s-1)\pi^{-s/2}\Gamma(s/2)\zeta(s),\qquad
\mathcal M F_h=g/h.
\end{equation}
Here $Z$ is the original nonempty set of actual nontrivial zeros,
closed under conjugation and $s\mapsto1-s$, with every complete
order $m_\rho$. This refers to the specified packet, not to the
existence of any off-line zero. The original Mellin convention is
$\mathcal M F(s)=\int_0^\infty F(x)x^{s-1}\,dx$ and $D=-x\partial_x$.
For its actual tensor degree $k\geq1$, put
\begin{equation}\tag{NS2}\label{NS2}
\begin{gathered}
c=k/2,\qquad S=c+iu,\qquad
w_h(t)=\frac{|(g/h)(1/2+it)|^2}{2\pi},\qquad m=w_h^{*k},\\
\mu_h=\int_\R w_h(t)\,dt,\qquad
\int_\R m(u)\,du=\mu_h^k=:\mathfrak m,\\
\mathcal A P=P(D_1+\cdots+D_k)F_h^{\otimes k},\qquad
\langle\mathcal A P,\mathcal A Q\rangle
 =\int_\R\overline{P(c+iu)}Q(c+iu)m(u)\,du .
\end{gathered}
\end{equation}
The last equality follows by Mellin Parseval on each original
$dx_j$ axis and by the convolution definition. In particular
$\|\mathcal A1\|^2=\mathfrak m$, not one.
The entire quotient $g/h$ is nonzero, so its isolated zeros make
$w_h$ positive almost everywhere. It is even by conjugation.
The vertical Gamma bound and the alternating-series integral for
$\zeta$ give a polynomial times $e^{-\pi|t|/2}$ for $w_h$ outside
a compact interval; the removable values of $g/h$ are bounded on
that interval. Thus all its moments exist, as do those of $m$.
Convolution preserves evenness and positivity almost everywhere.
Every nonzero polynomial therefore has positive squared norm.

Let $Q_n$ be the real monic orthogonal polynomials for this very
measure, with full norms
\begin{equation}\tag{NS3}\label{NS3}
\begin{gathered}
\omega_n=\int_\R Q_n(u)^2m(u)\,du>0,\qquad
\omega_0=\mathfrak m,\qquad a_n=\sqrt{\omega_{n+1}/\omega_n},
\quad a_{-1}=0,\\
\phi_n=Q_n/\sqrt{\omega_n},\qquad
p_n(S)=i^nQ_n((S-c)/i),\\
u\phi_n=a_{n-1}\phi_{n-1}+a_n\phi_{n+1},
\qquad Q_n(-u)=(-1)^nQ_n(u).
\end{gathered}
\end{equation}
Gram--Schmidt supplies the unique monic $Q_n$. Uniqueness after
$u\mapsto-u$ proves parity. Pairing $uQ_n$ with earlier
polynomials kills all terms below $Q_{n-1}$; parity kills the
$Q_n$ term; the leading coefficient and the $Q_{n-1}$ pairing
give the displayed recurrence. This proves the formulas without
substituting a comparison measure.

The exact original cyclic annihilator is
\begin{equation}\tag{NS4}\label{NS4}
\begin{gathered}
\chi(S)=\prod_\kappa(S-\kappa)^{e_\kappa},\qquad
e_\kappa=\max_{\rho_1+\cdots+\rho_k=\kappa}
 \left(1+\sum_{j=1}^k(m_{\rho_j}-1)\right),\\
q=\deg\chi,\qquad C=\C[S]/(\chi),\qquad
M=\frac{M_S-cI}{i},\qquad
\chi_u(z)=i^{-q}\chi(c+iz)=\det(zI-M).
\end{gathered}
\end{equation}
To prove the exponent formula, the local factor of
$(\C[s]/h)^{\otimes k}$ at $(\rho_1,\ldots,\rho_k)$ is
\[
\C[z_1,\ldots,z_k]/
       (z_1^{m_{\rho_1}},\ldots,z_k^{m_{\rho_k}}).
\]
Multiplication by $\sum s_j$ is $\kappa+\sum z_j$.
The last nonzero power of $\sum z_j$ has coefficient
$(\sum(m_{\rho_j}-1))!/\prod(m_{\rho_j}-1)!$ on
$\prod z_j^{m_{\rho_j}-1}$, which is nonzero.
The largest index over equal sums is precisely $e_\kappa$.
Consequently the substitution $\beta[P]=P(\sum s_j)$ is
injective on $C$. Multiplication by $S$ is its companion
operator in the original increasing-power remainder basis;
its minimal and characteristic polynomial are both $\chi$.
The substitution $S=c+iu$ proves all factors in $\chi_u$.
Because the original packet is conjugation and reflection
stable, $\chi_u$ has real coefficients and
$\chi_u(-z)=(-1)^q\chi_u(z)$.

All matrices through \eqref{NS21} use the original
$1,S,\ldots,S^{q-1}$ remainder coordinates, unless a different
map is explicitly displayed. Fix the actual degree $L\geq q-1$,
and write $b_n=[p_n]_\chi$. Define
\begin{equation}\tag{NS5}\label{NS5}
\begin{gathered}
E=\left[\frac{i^{-n}b_n}{\sqrt{\omega_n}}\right]_{n=0}^L,
\qquad K=EE^*=\sum_{n=0}^L\frac{b_nb_n^*}{\omega_n},
\qquad G=K^{-1},\\
R=E^*G,\qquad P=RE,\qquad Z=I-P.
\end{gathered}
\end{equation}
The first $q$ monic columns span $C$, so $E$ is onto and $G$
is positive definite. Direct multiplication gives $ER=I$,
$P=P^*=P^2$, and $R^*R=G$.
The kernel of $E$ is exactly the degree-$L$ polynomial
relations. Every lift of $v\in C$ is $Rv+z$ with $Ez=0$;
orthogonality gives its squared norm $v^*Gv+\|z\|^2$.
Thus $G$ is the attained original quotient metric.
Let $J_L$ be the real symmetric $(L+1)$-square Jacobi matrix
with entries $a_n$ above and below its diagonal. The next
section concerns the \emph{attained compression}
$EJ_LR$, not the uncompressed algebraic quotient $M$.

\section{The original outgoing edge and induced selfadjoint operator}

\begin{theorem}[Exact finite receiver]\label{thm:receiver}
In the original coordinates of \eqref{NS5},
\begin{equation}\tag{NS6}\label{NS6}
A:=EJ_LR=M+r\ell,\qquad
r=\frac{i}{\omega_L}b_{L+1},\qquad
\ell=b_L^*G,\qquad A^*G=GA .
\end{equation}
Here $\ell$ is the displayed row, not an unspecified
conjugate functional. In particular every eigenvalue of $A$
is real and $A$ is diagonalizable. Its eigenvectors can be
chosen orthonormal in $G$.
\end{theorem}
\begin{proof}
The exact rectangular matrix of multiplication by $u$ has
the additional outgoing term $a_Le_{L+1}e_L^*$.
Reduction of that rectangular matrix commutes with
multiplication by $u$ in $C$:
$E_{L+1}J_{L+1,L}=ME$. Removing the outgoing term and
composing with $R$ gives
\[
EJ_LR=M-a_L\frac{i^{-(L+1)}b_{L+1}}{\sqrt{\omega_{L+1}}}
              \frac{i^Lb_L^*G}{\sqrt{\omega_L}}.
\]
Since $i^{-(L+1)}i^L=-i$ and
$a_L/\sqrt{\omega_{L+1}\omega_L}=1/\omega_L$,
this is \eqref{NS6}. Further,
$GA=R^*J_LR$ is Hermitian. The positive square root of $G$
conjugates $A$ to the Hermitian matrix $G^{1/2}AG^{-1/2}$.
The finite Hermitian spectral theorem proves the assertions.
\end{proof}

Let $\Gamma:C\to C$ be $[P(S)]\mapsto[P(k-S)]$.
Its original remainder matrix is explicitly
\begin{equation}\tag{NS7}\label{NS7}
\Gamma_{rj}=\binom jr k^{j-r}(-1)^r\quad(0\leq r\leq j<q).
\end{equation}
Reflection stability of $\chi$ proves it is well-defined.
It satisfies $\Gamma^2=I$, $\Gamma M\Gamma=-M$ and
$\Gamma b_n=(-1)^nb_n$.
In native coefficient space let $\gamma=\operatorname{diag}((-1)^n)$.
Then $\Gamma E=E\gamma$, whence
$\Gamma K\Gamma^*=K$, $\Gamma^*G\Gamma=G$ and
$R\Gamma=\gamma R$. As $J_L\gamma=-\gamma J_L$, it follows that
\begin{equation}\tag{NS8}\label{NS8}
\Gamma A\Gamma=-A,\qquad
\Gamma(r\ell)\Gamma=-r\ell,\qquad
\ell r=\frac{i}{\omega_L}b_L^*Gb_{L+1}=0,\qquad
(r\ell)^2=0 .
\end{equation}
For the zero pairing, insert $\Gamma^*G\Gamma=G$ between the
opposite parity vectors to get the negative of the same scalar.
Thus the rank-one correction is square-zero, including the
case where it is zero. This does \emph{not} make $A$ and $M$
isospectral: mixed products of $M$ with $r\ell$ need not vanish.

For later exact estimates, put
\[
\alpha=b_L^*Gb_L,\qquad \delta=b_{L+1}^*Gb_{L+1}.
\]
The operator norm in this same metric is
\begin{equation}\tag{NS9}\label{NS9}
\|A-M\|_G=\frac{\sqrt{\alpha\delta}}{\omega_L}.
\end{equation}
Indeed $G^{1/2}r\ell G^{-1/2}$ is an outer product, of norm
$\|G^{1/2}r\|\|\ell G^{-1/2}\|$; inserting the two columns
proves the formula with the mass factors retained.
For every original eigenvalue $\alpha_0$ of $M$, choose a
nonzero eigenvector $v$, which exists even for a multiple root.
In a $G$-orthonormal eigenbasis for $A$,
$\|(A-\alpha_0)v\|_G\geq
\operatorname{dist}(\alpha_0,\operatorname{spec}A)\|v\|_G$.
The equality $(A-\alpha_0)v=(A-M)v$ therefore proves
\begin{equation}\tag{NS10}\label{NS10}
|\Im\alpha_0|
\leq\operatorname{dist}(\alpha_0,\operatorname{spec}A)
\leq\frac{\sqrt{\alpha\delta}}{\omega_L}.
\end{equation}
For the original root $\kappa=c+i\alpha_0$ this reads
$|\Re\kappa-c|\leq\sqrt{\alpha\delta}/\omega_L$.
It is a finite exact inequality, not a claim that its right
side tends to zero as any parameter grows.

\section{The full secular polynomial, including every local pole}

Define the actual rational function and polynomials
\begin{equation}\tag{NS11}\label{NS11}
\begin{gathered}
f(z)=\ell(zI-M)^{-1}r,\qquad
n(z)=\ell\,\operatorname{adj}(zI-M)\,r,\\
d(z)=\det(zI-A)
=\chi_u(z)-n(z)
=\chi_u(z)(1-f(z)).
\end{gathered}
\end{equation}
For $z$ outside the spectrum of $M$, factor $zI-A$ as
$(zI-M)(I-(zI-M)^{-1}r\ell)$. The determinant of $I-v\ell$
is $1-\ell v$: choose a basis beginning with $v$ if it is
nonzero, and read its diagonal. This proves the rational
identity. Multiplication by $\chi_u$ and the adjugate
identity extend it to a polynomial identity at every root
of $\chi_u$. It is crucial that no value is assigned to
$f$ at a pole by omitting that pole.

Parity gives a stronger degree statement:
\begin{equation}\tag{NS12}\label{NS12}
f(-z)=f(z),\qquad
n(-z)=(-1)^qn(z),\qquad
d(-z)=(-1)^qd(z),\qquad
\deg n\leq q-2 .
\end{equation}
The function $f$ is even because
$(-zI-M)^{-1}=-\Gamma(zI-M)^{-1}\Gamma$,
$\ell\Gamma=(-1)^L\ell$, and $\Gamma r=(-1)^{L+1}r$.
Also
$f(z)=\sum_{j\geq0}\ell M^jr\,z^{-j-1}$ near infinity.
The coefficient at $z^{-1}$ is $\ell r=0$, proving the degree
bound. More explicitly,
\begin{equation}\tag{NS13}\label{NS13}
\ell M^{2j}r=0,\qquad
f(z)=\sum_{j\geq0}\frac{\ell M^{2j+1}r}{z^{2j+2}}
\quad(|z|>\|M\|).
\end{equation}
Insert the parity involution in the scalar to verify each
vanishing coefficient. For $q=1$, the degree bound means
$n=0$, and the correction itself is zero because $r\ell$
is a scalar with trace zero.

Here are all local coefficients, with repeated roots retained.
Let $\alpha$ range over distinct roots of $\chi_u$ and let
$e_\alpha$ be its full multiplicity. Set
$\chi_{\alpha}(z)=\chi_u(z)/(z-\alpha)^{e_\alpha}$ and define
the exact Chinese-remainder polynomial
\begin{equation}\tag{NS14}\label{NS14}
e_\alpha(z)=\chi_\alpha(z)
\sum_{a=0}^{e_\alpha-1}
\frac{(z-\alpha)^a}{a!}
\left(\frac{d^a}{dz^a}\frac1{\chi_\alpha(z)}\right)_{z=\alpha},
\qquad \Pi_\alpha=e_\alpha(M).
\end{equation}
Its residues are $1$ modulo $(z-\alpha)^{e_\alpha}$ and
$0$ modulo every other primary factor. Therefore
$\Pi_\alpha\Pi_\beta=\delta_{\alpha\beta}\Pi_\alpha$,
$\sum_\alpha\Pi_\alpha=I$, and
$(M-\alpha I)^{e_\alpha}\Pi_\alpha=0$.
The resolvent is exactly the finite sum
\begin{equation}\tag{NS15}\label{NS15}
\begin{gathered}
(zI-M)^{-1}
 =\sum_\alpha\sum_{j=0}^{e_\alpha-1}
       \frac{(M-\alpha I)^j\Pi_\alpha}{(z-\alpha)^{j+1}},\\
f(z)=\sum_\alpha\sum_{j=0}^{e_\alpha-1}
       \frac{c_{\alpha,j}}{(z-\alpha)^{j+1}},\qquad
c_{\alpha,j}=\ell(M-\alpha I)^j\Pi_\alpha r .
\end{gathered}
\end{equation}
Multiply the first sum by $zI-M$ on each primary space:
the finite geometric series cancels, leaving $\Pi_\alpha$.
This proves the formula and all its coefficients.
Let $p_\alpha$ be zero when all $c_{\alpha,j}$ vanish,
and otherwise one plus the largest $j$ with
$c_{\alpha,j}\ne0$.

\begin{theorem}[Exact cancellation and the invisible factor]
\label{thm:cancellation}
Write the fully reduced rational expression as
\begin{equation}\tag{NS16}\label{NS16}
f(z)=\frac{P_f(z)}{D_f(z)},\qquad
D_f(z)=\prod_\alpha(z-\alpha)^{p_\alpha},\qquad
\gcd(P_f,D_f)=1,
\end{equation}
where $D_f$ is monic. When $f=0$ set $P_f=0,D_f=1$.
Then
\begin{equation}\tag{NS17}\label{NS17}
d(z)=I_f(z)H_f(z),\qquad
I_f(z)=\frac{\chi_u(z)}{D_f(z)},\qquad
H_f(z)=D_f(z)-P_f(z).
\end{equation}
The factor $H_f$ is monic, has only real roots, and every
one of its roots is simple. The factor $I_f$ also has only
real roots, but its roots and those of $H_f$ need not be
disjoint. No invisible factor is discarded.

At an original root $\alpha$ of multiplicity $e_\alpha$,
the exact multiplicity in $d$ is
\begin{equation}\tag{NS18}\label{NS18}
\ord_\alpha d=
\begin{cases}
e_\alpha-p_\alpha,&p_\alpha>0,\\
e_\alpha+\ord_\alpha(1-f),&p_\alpha=0,
\end{cases}
\end{equation}
where in the second case $\ord_\alpha(1-f)$ is either
zero or one. In particular every nonreal original root
has $p_\alpha=e_\alpha$. A real root which is absent from
$d$ also has $p_\alpha=e_\alpha$.
\end{theorem}
\begin{proof}
The finite partial fractions prove \eqref{NS16}, including
the precise pole exponents. The numerator has lower degree
than the denominator, since $f$ vanishes at infinity.
Equation \eqref{NS11} now proves \eqref{NS17}.

The resolvent identity, proved by multiplying both sides
by $zI-M-r\ell$, is
\begin{equation}\tag{NS19}\label{NS19}
(zI-A)^{-1}
=(zI-M)^{-1}
 +\frac{(zI-M)^{-1}r\ell(zI-M)^{-1}}{1-f(z)}.
\end{equation}
It holds off both finite spectra and then as a rational
identity. Applying $\ell$ and $r$ gives
\[
\ell(zI-A)^{-1}r=\frac{f(z)}{1-f(z)}
                 =\frac{P_f(z)}{D_f(z)-P_f(z)}.
\]
Because $\gcd(P_f,D_f)=1$, the last numerator and
denominator are coprime. The $G$-selfadjoint spectral
decomposition of $A$ writes its resolvent as
$\sum_t\mathsf P_t/(z-t)$ over its distinct real
eigenvalues, with $G$-orthogonal projections $\mathsf P_t$.
Every scalar resolvent thus has at most simple real poles.
The reduced denominator $H_f$ must therefore have simple
real roots. When $f=0$, $H_f=1$ and this assertion is empty.
All roots of $d$ are real, hence its factor $I_f$ has
only real roots as well.

If $p_\alpha>0$, the leading pole coefficient of $f$
is nonzero, so $1-f$ has that same pole order; this proves
the first case of \eqref{NS18}. If $p_\alpha=0$,
$D_f(\alpha)\ne0$ and $1-f=H_f/D_f$ locally; simplicity
of $H_f$ proves the second case and its two alternatives.
Since $d$ has no nonreal root, its nonnegative local
order must be zero at each nonreal $\alpha$, forcing
the full cancellation $p_\alpha=e_\alpha$.
\end{proof}

There is a useful restriction on all surviving multiplicities.
Since $M$ is companion, $\dim\ker(M-\alpha I)=1$ at any
original root. On $\ker(A-\alpha I)$ the map $v\mapsto\ell v$
has a kernel contained in $\ker(M-\alpha I)$, and its image
has dimension at most one. Therefore
\begin{equation}\tag{NS20}\label{NS20}
\dim\ker(A-\alpha I)\leq2
\quad(\alpha\in\operatorname{spec}M),\qquad
\dim\ker(A-tI)\leq1
\quad(t\notin\operatorname{spec}M).
\end{equation}
Diagonalizability makes these the algebraic multiplicities.
For example, if $p_\alpha>0$ then $e_\alpha-p_\alpha\leq2$;
if $p_\alpha=0$ then
$e_\alpha+\ord_\alpha(1-f)\leq2$.
This conclusion includes shared active/invisible roots
without assuming their absence.

Parity pairs the local poles, including their signs:
\begin{equation}\tag{NS21}\label{NS21}
c_{-\alpha,j}=(-1)^{j+1}c_{\alpha,j},\qquad
p_{-\alpha}=p_\alpha,\qquad
c_{0,j}=0\quad\hbox{when }j+1\hbox{ is odd}.
\end{equation}
Expand $f(-z)=f(z)$ near $z=-\alpha$ to obtain these
identities. In particular $\deg D_f$ is even: nonzero
poles occur in equal opposite pairs and a pole at zero
has even order. Thus $D_f$ and $P_f$ are even, $H_f$
is even, and its nonzero simple real roots occur in
opposite pairs. Since an even polynomial has even
order at zero, $H_f(0)\ne0$. Any zero eigenvalue of
$A$ belongs entirely to $I_f$. These facts also show
directly why dropping a repeated original pole or the
invisible factor changes the answer.

\section{Original S coordinates, actual eigenvectors and physical jets}

There is no unannounced identification of $u$ with $S$.
Define the actual attained scaling operator
\begin{equation}\tag{NS22}\label{NS22}
B=cI+iA=M_S-\frac{b_{L+1}b_L^*G}{\omega_L}.
\end{equation}
It is $G$-normal, with $B^{\dagger_G}=cI-iA=kI-B$.
Its characteristic polynomial is exactly
\begin{equation}\tag{NS23}\label{NS23}
\begin{split}
d_S(s)&=\det(sI-B)=i^qd((s-c)/i)\\
 &=\chi(s)
  +\frac1{\omega_L}b_L^*G\,
               \operatorname{adj}(sI-M_S)b_{L+1}.
\end{split}
\end{equation}
The first equality takes the determinant of
$sI-B=i((s-c)I/i-A)$, including $i^q$.
The second follows from the rank-one determinant identity
applied to $sI-M_S+b_{L+1}\ell/\omega_L$.
Thus every root of $d_S$ is on $\Re s=c$.
This statement concerns $d_S$, not the different polynomial
$\chi$. Equation \eqref{NS23} is their exact relationship.

All eigenvectors can be recovered without a singular
evaluation of \eqref{NS19}. At a root $t$ of $d$ with
$\chi_u(t)\ne0$, the vector
\begin{equation}\tag{NS24}\label{NS24}
v_t=(tI-M)^{-1}r,\qquad
(tI-A)v_t=0,\qquad \ell v_t=f(t)=1
\end{equation}
is nonzero and spans the eigenspace, by \eqref{NS20}.
Its actual projector is
$v_t(v_t^*Gv_t)^{-1}v_t^*G$.
At a shared old root $t$, the complete nonsingular
system is instead
\begin{equation}\tag{NS25}\label{NS25}
\begin{pmatrix}tI-M&-r\\ \ell&-1\end{pmatrix}
\binom v\tau=0
\end{equation}
with scalar $\tau$; explicitly
$(tI-M)v=r\tau$, $\ell v=\tau$.
The map $(v,\tau)\mapsto v$ is a bijection from this
kernel onto $\ker(tI-A)$, with inverse
$v\mapsto(v,\ell v)$. It retains invisible eigenvectors
with $\ell v=0$ and does not insert a false inverse
at the old pole. If $V_t$ is any matrix whose columns
are a basis of the resulting eigenspace, then
\begin{equation}\tag{NS26}\label{NS26}
\mathsf P_t=V_t(V_t^*GV_t)^{-1}V_t^*G,\quad
\sum_t\mathsf P_t=I,\quad
A=\sum_t t\mathsf P_t,\quad
(zI-A)^{-1}=\sum_t\frac{\mathsf P_t}{z-t}.
\end{equation}
The Gram $V_t^*GV_t$ is positive definite because the
columns are independent. The expression is the
$G$-orthogonal projection onto their span. Distinct
eigenspaces of a $G$-selfadjoint operator are
$G$-orthogonal: taking the pairing in the two orders
gives $(t-t')\langle v',v\rangle_G=0$.
Completeness follows from the Hermitian conjugate used
in the proof of \eqref{NS6}. This proves all four formulas.

For precision, the original $u$-power remainder map is
\begin{equation}\tag{NS27}\label{NS27}
T_{rj}=\binom jr c^{j-r}i^r\quad(0\leq r\leq j<q),\qquad
[P(S)]_\chi\longmapsto[P(c+iu)]_{\chi_u}=T[P].
\end{equation}
It is triangular with nonzero diagonal $i^j$.
In those coordinates
$M_u^{\rm pow}=TMT^{-1}$,
$G_u=(T^{-1})^*GT^{-1}$ and
$d_n=[Q_n]_{\chi_u}=i^{-n}Tb_n$. Hence
\begin{equation}\tag{NS28}\label{NS28}
TAT^{-1}
=M_u^{\rm pow}-\frac{d_{L+1}d_L^*G_u}{\omega_L},\qquad
G_u=\left(\sum_{n=0}^L
\frac{d_nd_n^*}{\omega_n}\right)^{-1}.
\end{equation}
This proves the real-coordinate sign and both coordinate
metrics. It does not replace the original $S$ coordinates.

The physical injection retains the full local unit.
Put $E_h=\C[s]/h$, $\upsilon=j_h(g/h)$, and
\begin{equation}\tag{NS29}\label{NS29}
\eta:C\longrightarrow E_h^{\otimes k},\qquad
\eta[P]=\upsilon^{\otimes k}P(s_1+\cdots+s_k).
\end{equation}
At the root $\rho$, its complete unit coefficient is
\begin{equation}\tag{NS30}\label{NS30}
\begin{split}
u_{\rho,r}
={}&\sum_{a=0}^r
\frac{g^{(m_\rho+a)}(\rho)}{(m_\rho+a)!}
\sum_{\substack{\ell_\sigma\geq0\ (\sigma\ne\rho)\\
                 \sum_{\sigma\ne\rho}\ell_\sigma=r-a}}
\prod_{\sigma\ne\rho}
\frac{(-1)^{\ell_\sigma}
 \binom{m_\sigma+\ell_\sigma-1}{\ell_\sigma}}
 {(\rho-\sigma)^{m_\sigma+\ell_\sigma}},
\qquad 0\leq r<m_\rho .
\end{split}
\end{equation}
Expand $g(\rho+z)/z^{m_\rho}$ and each factor
$(\rho-\sigma+z)^{-m_\sigma}$ to verify the formula.
In particular
$u_{\rho,0}=g^{(m_\rho)}(\rho)/
(m_\rho!\prod_{\sigma\ne\rho}(\rho-\sigma)^{m_\sigma})\ne0$.
Thus multiplication by this full jet is invertible.
Together with the injectivity of $\beta$ proved in
\eqref{NS4}, this proves injectivity of $\eta$.
For each original local multi-index the map is
\begin{equation}\tag{NS31}\label{NS31}
(\eta[P])_{\boldsymbol\rho,\boldsymbol j}
=\sum_{0\leq\nu_a\leq j_a}
 \left(\prod_{a=1}^k u_{\rho_a,j_a-\nu_a}\right)
 \frac{P^{(|\boldsymbol\nu|)}(\rho_1+\cdots+\rho_k)}
      {\prod_{a=1}^k\nu_a!}.
\end{equation}
Taylor's formula for $P(\sum\rho_a+\sum z_a)$ gives its
coefficient $P^{(|\boldsymbol\nu|)}/\prod\nu_a!$;
multiplication by the unit jets proves \eqref{NS31}.
All phases, multiplicities, sum collisions and unit
derivatives are present.

Give $\eta(C)$ the actual transported metric
\[
\langle\eta v,\eta w\rangle_\eta=v^*Gw.
\]
Then
\begin{equation}\tag{NS32}\label{NS32}
A_\eta=\eta A\eta^{-1},\quad
B_\eta=\eta B\eta^{-1},\quad
\mathsf P_{t,\eta}=\eta\mathsf P_t\eta^{-1}
\end{equation}
are exact operators on this image, with the same
determinants and the same selfadjoint/normal statements.
The inverse is used only on $\eta(C)$.
For a coordinate basis of this image, these formulas
mean the usual congruence $(\eta^{-1})^*G\eta^{-1}$,
not a Euclidean Gram on a different jet space.
The physical representative of $\eta v$ is
$\mathcal A\sum_{n=0}^L (Rv)_n i^{-n}p_n/\sqrt{\omega_n}$.
Its norm is $v^*Gv$, its jet is $\eta v$, and replacing
$v$ by $Av$ or $Bv$ gives the corresponding attained
physical receiver. These equalities follow from
$ER=I$, \eqref{NS2}, and the definition of $E$.

\section{What the human-source calculation supplies}

In \cite{CCM}, Lemma \texttt{key}, the original finite
data are a diagonal $D$ with entries $-N,\ldots,N$ and
a positive semidefinite Weil matrix $T$ with a
one-dimensional even radical $\C\xi$.
The authors construct $D'=D-D\xi\,\eta^*$ and descend
through that radical, obtaining a selfadjoint $D''$ in
the \emph{induced} $T$ metric. Their determinant is
\[
\det(D''-s)=\det(D-s)\sum_{j=-N}^N\frac{\xi_j}{j-s}.
\]
Its proof uses exactly the finite rank-one determinant
identity, while the quotient by the radical removes
the additional eigenvalue zero of $D'$.
The present \eqref{NS6}--\eqref{NS19} apply that algebra
to a different, explicitly identified receiver: the
full arithmetic cyclic quotient with positive attained
metric $G$, outgoing native polynomial column, and
possibly repeated nonreal original poles.
Equations \eqref{NS14}--\eqref{NS18} calculate those
poles rather than treating the diagonal simple-pole
formula as if it already included them.

The original source's Proposition \texttt{pertscal}
and Theorem \texttt{finmain} retain the periodic
scaling factor $2\pi/(2\log\lambda)$ and spectral-cut
phase $-i\lambda^{-iz}$ in its regularized determinant.
Neither factor belongs to the finite polynomial
determinant \eqref{NS11}; the exact original scaling
map here is instead \eqref{NS22}--\eqref{NS23}.
The source's Sections \texttt{outlook} and
\texttt{missing} explicitly leave the minimum
eigenvector/prolate approximation and the limiting
zero argument unproved. Nothing in this finite note
supplies those limiting estimates.

The completed finite conclusion is stronger than a
real-zero assertion for an unrelated matrix: the
original $\chi$, attained metric, outgoing relation
column and complete physical jet injection are
connected by an exact determinant and resolvent
identity. Every nonreal original factor is cancelled
by the full pole order of the scalar secular term
when forming $d=\chi_u(1-f)$, and every
surviving original factor remains explicitly in
\eqref{NS17}. Neither cancellation nor finite reality
changes the original zeros of $\chi$.

\section{An exact finite example with nonreal original poles}

The following algebra check is not a fabricated arithmetic
packet. It uses the explicitly different measure
$(3/2)\,1_{[-1,1]}(u)\,du$ of mass $3$ and the test quotient
$\chi_u(z)=z^2+1$, at $L=1$. Its monic polynomials are
$Q_0=1$, $Q_1=u$, $Q_2=u^2-1/3$ with
$\omega_0=3$, $\omega_1=1$. In the $1,u$ remainder basis,
\begin{equation}\tag{NS33}\label{NS33}
M=\begin{pmatrix}0&-1\\1&0\end{pmatrix},\quad
G=\begin{pmatrix}3&0\\0&1\end{pmatrix},\quad
A-M=\begin{pmatrix}0&4/3\\0&0\end{pmatrix},\quad
A=\begin{pmatrix}0&1/3\\1&0\end{pmatrix}.
\end{equation}
Direct multiplication verifies $A^*G=GA$ and
$(A-M)^2=0$. Nevertheless
\begin{equation}\tag{NS34}\label{NS34}
f(z)=\frac{4/3}{z^2+1},\qquad
d(z)=z^2-\frac13,\qquad
D_f=z^2+1,\quad P_f=4/3,\quad I_f=1.
\end{equation}
The poles at $i$ and $-i$ are both full-order and
cancelled, while the new zeros are $\pm1/\sqrt3$.
This exact example falsifies the inference that a
square-zero correction leaves eigenvalues unchanged,
or that reality of $d$ by itself proves reality of
$\chi_u$. It also checks the mass and correction sign.
The executable verification accompanying this note
uses additional exact positive-measure models for
repeated poles and invisible factors. Those tests
verify finite algebra; they do not numerically
evaluate the arithmetic measure in \eqref{NS2}.

\section{The same original rank-two arithmetic allowance}

The outgoing secular correction also recovers the
original Hermitian control, without changing its
metric. Write $Q=A-M=r\ell$ and retain
$X^{\dagger_G}=G^{-1}X^*G$. Since $A^{\dagger_G}=A$,
$M=A-Q$ gives
\begin{equation}\tag{NS35}\label{NS35}
\begin{split}
\mathcal C
&=M_S+M_S^{\dagger_G}-kI
=i(M-M^{\dagger_G})\\
&=-i(Q-Q^{\dagger_G})
=\frac{b_{L+1}b_L^*+b_Lb_{L+1}^*}{\omega_L}\,G .
\end{split}
\end{equation}
Thus $\mathcal C^{\dagger_G}=\mathcal C$; equivalently
$G\mathcal C$ is an ordinary Hermitian matrix.
Put $x=G^{1/2}b_{L+1}$ and $y=G^{1/2}b_L$.
Parity gives $y^*x=0$. Conjugating $\mathcal C$ by
$G^{1/2}$ gives $(xy^*+yx^*)/\omega_L$.
If either vector is zero this operator is zero.
Otherwise its action in the orthonormal pair
$x/\|x\|,y/\|y\|$ is
$\left(\begin{smallmatrix}0&\|x\|\|y\|/\omega_L\\
\|x\|\|y\|/\omega_L&0\end{smallmatrix}\right)$;
on the orthogonal complement it vanishes. Consequently
its two possibly nonzero eigenvalues are
\begin{equation}\tag{NS36}\label{NS36}
\pm\epsilon_L,\qquad
\epsilon_L=
\frac{\sqrt{(b_L^*Gb_L)(b_{L+1}^*Gb_{L+1})
                 -|b_L^*Gb_{L+1}|^2}}{\omega_L}
=\|Q\|_G .
\end{equation}
The first expression retains the complete two-vector
Gram determinant; the middle pairing is zero here
by the proved packet parity, not by deleting a term.
For any original eigenvector $M_Sv=\kappa v$,
\[
\frac{v^*G\mathcal C v}{v^*Gv}=2\Re\kappa-k.
\]
The spectral bounds just proved yield the sharper
original-coordinate finite estimate
\begin{equation}\tag{NS37}\label{NS37}
|2\Re\kappa-k|\leq\epsilon_L
\qquad(\kappa\in\operatorname{spec}M_S).
\end{equation}
Thus the scalar secular replacement and the original
rank-two arithmetic allowance receive exactly the
same outgoing column. This strengthens the strip
consequence of \eqref{NS10} by the factor two
visible in the original scaling generator; the
distance-to-spectrum comparison in that equation
remains a separate bound. Neither
finite bound gives a decay estimate for $\epsilon_L$.

\begin{thebibliography}{9}\raggedright
\bibitem{CCM} Alain Connes, Caterina Consani and Henri Moscovici,
\emph{Zeta Spectral Triples},
\href{https://arxiv.org/abs/2511.22755v1}{arXiv:2511.22755v1},
27 November 2025. Original author source
\texttt{mc2arXiv.tex}. Read for this receiver:
Section 5 from truncated matrices through the complete
proof of Theorem \texttt{finmain}, including Lemma
\texttt{key}, equation \texttt{eq:detDp}, Proposition
\texttt{pertscal}, the regularized determinant proof;
Sections \texttt{outlook} and \texttt{missing}.
Source SHA-256:
\texttt{cd02adff07e89dcd343cf0dfc06a14d7190cef43e1978e20706028dda92fa71b}.
No PDF was read. Books cited within the outlook were
not independently read for this receiver.
\bibitem{Native} Split-Zero research programme,
\emph{The native arithmetic prolate action, its Lie
defects, and the exact polynomial-quotient receiver},
20 September 2026, result SZ-20260920-031,
\texttt{NATIVE\_PROLATE\_LIE\_DEFECT.tex},
especially equations \texttt{finitequotient},
\texttt{Jquotient}, \texttt{unitreceiver} and
\texttt{etajets}. Source SHA-256:
\texttt{7f8fa549c99b725a8641c43419e561e8bf0ac90c5a37e342d797a80651d76a3f}.
The finite quotient, outgoing correction and unit
formulas used here have been proved again above.
\end{thebibliography}
\end{document}
